% ======================================================================
% CAPITULO 14 — OpenCode Ecosystem como Laboratorio GAT
% ======================================================================
\chapter{OpenCode Ecosystem como Laboratorio GAT}

\section{Por Que um Laboratorio de IA para GAT?}

A Geometric Arbitrage Theory e matematicamente elegante, mas \textbf{validar}
seus resultados empiricamente requer ferramentas que vao alem do calculo
manual. O \textbf{OpenCode Ecosystem}\footnote{Claro, M. \textit{OpenCode
Ecosystem v4.6.1: Arquitetura Multiagente Evolutiva}. GitHub, 2026.
\url{https://github.com/MarceloClaro/OpenCode\_Ecosystem}. \textbf{Resenha:}
Plataforma multi-agente de pesquisa cientifica assistida por IA. Integra 79
agentes especializados, 38 servidores MCP, 104 skills e 7 verificadores
simbolicos Cora-Debate (V1--V7). Projetada para transformar hipoteses
cientificas em resultados verificaveis com rastreabilidade completa
(principio de integridade: INTEGRIDADE.md).} e uma plataforma de 79 agentes
especializados, 38 servidores MCP e 7 verificadores simbolicos que oferece
exatamente o ambiente necessario para:

\begin{enumerate}[nosep]
  \item \textbf{Simular mercados} com precos estocasticos e testar a presenca
        de arbitragem
  \item \textbf{Calcular curvatura} do fibrado do mercado automaticamente
  \item \textbf{Verificar} que condicoes produzem $R=0$ e quais violam NFLVR
  \item \textbf{Documentar} cada passo com rastreabilidade completa (SDD+TDD)
  \item \textbf{Evoluir} o raciocinio cientifico a cada ciclo de experimentacao
\end{enumerate}

\subsection{O Que E o OpenCode Ecosystem?}

Para quem nunca usou, o OpenCode e uma plataforma de \textbf{IA multi-agente}
--- em vez de um unico modelo de linguagem respondendo perguntas, dezenas de
agentes especializados colaboram, debatem e verificam o trabalho uns dos
outros\footnote{Para uma introducao visual ao conceito de agentes de IA, veja
Wang et al. (2024). \textit{A Survey on Large Language Model based Autonomous
Agents}. Frontiers of Computer Science. DOI: 10.1007/s11704-024-40231-1.
\textbf{Resenha para leigos:} Este artigo explica como funciona um ``time''
de IAs que trabalham juntas --- cada uma com uma especialidade diferente
(matematica, programacao, revisao critica). E como ter uma banca de
professores, cada um checando uma parte diferente do seu trabalho.}.

\begin{center}
\fbox{\begin{minipage}{0.88\textwidth}
\textbf{Analogia para leigos:} Imagine que voce tem 79 estagiarios muito
inteligentes. Um busca artigos, outro escreve codigo, outro verifica as contas,
outro procura erros, outro sugere melhorias. Eles trabalham em paralelo e
depois debatem entre si. O resultado final passa por 7 verificadores
automaticos (os ``Cora-Debate'') antes de ser aprovado. Voce, como
pesquisador, define a pergunta e interpreta os resultados --- mas o trabalho
braçal de verificacao e feito pelo time de agentes.
\end{minipage}}
\end{center}

\section{Instalacao e Primeiro Experimento}

\subsection{Instalacao em 3 Passos}

\begin{lstlisting}[caption={Instalacao do OpenCode Ecosystem}, label=lst:install]
# Passo 1: Clonar o repositorio
git clone https://github.com/MarceloClaro/OpenCode_Ecosystem.git
cd OpenCode_Ecosystem

# Passo 2: Instalar dependencias
bun install
pip install -r requirements.txt

# Passo 3: Verificar instalacao
python -c "from core.container import Container; print('OK')"
\end{lstlisting}

\subsection{Experimento 1: Simular um Mercado de 2 Ativos}

Nosso primeiro experimento\footnote{Este experimento e uma simplificacao
didatica. Mercados reais tem milhares de ativos, custos de transacao,
restricoes regulatorias e latencia de execucao --- nenhum dos quais e
modelado aqui. O objetivo e ilustrar o conceito, nao construir uma estrategia
de trading. Conforme principio de integridade R-I3 (INTEGRIDADE.md):
 distinguimos o que e medido do que e projetado.} e criar um mercado
sintetico com 2 ativos e verificar se existe arbitragem:

\begin{lstlisting}[caption={Experimento 1: Mercado de 2 ativos com GBM}]
import numpy as np
from scipy.stats import norm

# Parametros do mercado
S0 = np.array([100.0, 100.0])  # precos iniciais
mu = np.array([0.05, 0.08])    # retornos esperados (diferentes!)
sigma = np.array([0.20, 0.25]) # volatilidades
r = 0.03                        # taxa livre de risco

# Simular 1 trajetoria (Euler-Maruyama)
T, N = 1.0, 252
dt = T / N
S = np.zeros((N+1, 2))
S[0] = S0
for i in range(N):
    dW = np.random.normal(0, np.sqrt(dt), 2)
    S[i+1] = S[i] * (1 + mu*dt + sigma*dW)

# Calcular precos descontados
discount = np.exp(-r * np.linspace(0, T, N+1))
S_hat = S * discount.reshape(-1, 1)

# Verificar: se mu != r, entao S_hat nao e martingal
# → existe arbitragem → curvatura R != 0
print(f"Retorno esperado anualizado: {mu}")
print(f"Taxa livre de risco: {r}")
print(f"Ha arbitragem? {'SIM' if any(mu != r) else 'NAO'}")
\end{lstlisting}

\textbf{Resultado esperado:} Como $\mu_2 = 0,08 \neq 0,03 = r$, o ativo 2
oferece retorno esperado acima da taxa livre de risco \textbf{na medida
fisica} $P$. Isso \textbf{nao} constitui arbitragem --- e apenas um premio
de risco. A arbitragem so existe se \textbf{nao for possivel} encontrar uma
medida $P^*$ sob a qual todos os retornos esperados (descontados) sejam iguais
a $r$.

\subsection{Experimento 2: Construir uma Arbitragem Real}

Para construir uma arbitragem, precisamos de \textbf{inconsistencia de precos}
--- dois ativos que deveriam ter o mesmo preco (porque sao equivalentes) mas
nao tem:

\begin{lstlisting}[caption={Experimento 2: Arbitragem por inconsistencia}]
# Dois ativos identicos precificados diferentemente
S1 = np.array([100, 102, 101, 103, 105])  # Ativo 1
S2 = np.array([100, 101, 100, 102, 104])  # Ativo 2 (deveria ser = S1)

# Estrategia: comprar o barato, vender o caro
diferenca = S1 - S2
posicao = -np.sign(diferenca)  # -1: vender S1, comprar S2
pnl = posicao[:-1] * (diferenca[1:] - diferenca[:-1])

print(f"Diferenca maxima de precos: {max(abs(diferenca))}")
print(f"PnL acumulado: {sum(pnl):.2f}")
print(f"Sharpe ratio: {np.mean(pnl)/np.std(pnl):.2f}")
\end{lstlisting}

\section{Usando os Verificadores Cora-Debate (V1-V7)}

O coracao da verificacao no OpenCode sao os 7 verificadores
Cora-Debate\footnote{O nome ``Cora'' vem de \textit{Cognitive ORchestrated
Argumentation} --- um sistema que simula o processo de revisao por pares
cientificos. Cada verificador e como um revisor especializado: um checa a
matematica (V1), outro a logica (V3), outro a estatistica (V4), e assim por
diante. Juntos, eles formam uma ``banca'' automatica que audita cada
afirmacao.}. Cada um audita uma dimensao diferente do raciocinio:

\begin{center}
\begin{tabular}{p{0.05\textwidth}p{0.20\textwidth}p{0.55\textwidth}}
\toprule
\textbf{V} & \textbf{Nome} & \textbf{O que verifica (explicacao para leigos)} \\
\midrule
V1 & Dimensional & As unidades de medida batem? (ex: nao se soma kg com metros) \\
V2 & Algebrico (SymPy) & As equacoes foram derivadas corretamente? \\
V3 & Contraexemplos & Existe algum caso onde a afirmacao e falsa? (principio de Popper) \\
V4 & Estatistico & Os p-valores e intervalos de confianca estao corretos? \\
V5 & Numerico & A precisao dos calculos e suficiente? Ha erros de arredondamento? \\
V6 & EDO/EDP & Equacoes diferenciais satisfazem condicoes de existencia? \\
V7 & Codigo-fonte & O codigo que gerou os resultados esta correto? Ha vulnerabilidades? \\
\bottomrule
\end{tabular}
\end{center}

\subsection{Exemplo: Validando o Calculo de Curvatura com V1-V7}

\begin{lstlisting}[caption={Validacao do calculo de curvatura com Cora-Debate}]
# Script: validar_curvatura_gat.py
# Executar: opencode run /validar_gat

import sympy as sp
import numpy as np

# Definir simbolos
r, t = sp.symbols('r t')
D1, D2 = sp.symbols('D1 D2', cls=sp.Function)
x1, x2 = sp.symbols('x1 x2')

# Calcular curvatura para mercado de 2 ativos
# R = dchi + [chi,chi]  (como G* e abeliano, [chi,chi] = 0)
# chi = (1/D^x) * sum(D^j dx^j) - r^x dt

D_x = x1*D1(t) + x2*D2(t)
r_x = (x1*D1(t)*r + x2*D2(t)*r) / D_x  # simplificado: mesma taxa r

# Calcular derivada exterior dchi...
# (implementacao completa no repositorio)

# Verificacao V1 (dimensional):
# [r] = 1/T, [D] = valor → [chi] = adimensional? ✓
# [R] = 1/(valor × tempo)? Verificar.

# Verificacao V2 (algebrico):
# R == 0 quando D1 e D2 satisfazem dD/D = r dt
# (movimento Browniano geometrico com drift = r)

# Verificacao V3 (contraexemplo):
# Se D1 = exp(r t + sigma W_t) e D2 = exp(r t - sigma W_t),
# entao R != 0 → existe arbitragem
\end{lstlisting}

\section{SDD+TDD: Desenvolvimento Verificado de Modelos GAT}

O ecossistema aplica \textbf{Spec-Driven Development} e \textbf{Test-Driven
Development} a modelos financeiros\footnote{Beck, K. \textit{Test-Driven
Development: By Example}. Addison-Wesley, 2003. ISBN: 978-0321146533.
\textbf{Resenha para leigos:} TDD e como construir uma casa comecando pelos
alicerces: voce primeiro escreve o que quer testar (``a parede deve aguentar
X quilos''), depois constroi a parede, e depois verifica se ela aguenta.
Na GAT, voce primeiro escreve ``a curvatura deve ser zero quando mu = r'',
depois implementa o codigo, e depois verifica se o teste passa.}:

\begin{enumerate}[nosep]
  \item \textbf{SPEC:} Escrever a especificacao do modelo (ex: ``mercado com
        $N$ ativos de Ito, conexao $\chi$, curvatura $R$ deve satisfazer
        $R=0 \Leftrightarrow \mu_j = r$ para todo $j$'')
  \item \textbf{TDD:} Escrever o teste \textbf{antes} do codigo:
        \texttt{assert calcular\_curvatura(mu=[r,r], sigma=[0.2,0.2]) == 0}
  \item \textbf{IMPLEMENTAR:} Codigo que calcula $R$ a partir de $\chi$
  \item \textbf{VERIFICAR:} Executar o teste. Se falhar (RED), corrigir o codigo
        ate passar (GREEN). Refatorar (REFACTOR).
  \item \textbf{AUDITAR:} Submeter ao Cora-Debate V1--V7 para verificacao
        simbolica.
\end{enumerate}

\subsection{Exemplo de SPEC para um Modelo GAT}

\begin{lstlisting}[caption={SPEC: Modelo GAT de 3 ativos com Ito}]
"""
SPEC-GAT-001: Mercado de 3 ativos com precos de Ito

ENTRADAS:
  - N=3 ativos
  - precos: dS_t^j/S_t^j = mu_j dt + sigma_j dW_t^j
  - correlacoes: d[W^i, W^j]_t = rho_ij dt
  - taxa livre de risco: r

SAIDAS:
  - R(x,t,g): 2-forma de curvatura
  - Hol(chi): grupo de holonomia

CRITERIOS DE ACEITACAO (CTs):
  CT-1: Se mu_j = r para todo j, entao R ≡ 0 (NFLVR)
  CT-2: Se mu_j != r para algum j, entao R != 0
  CT-3: dim(Hol(chi)) = numero de mu_j != r (estrat. independentes)
  CT-4: V1-V7: todos os verificadores passam
"""
\end{lstlisting}

\section{Integracao com o Pipeline Academico}

O OpenCode oferece um pipeline completo\footnote{O pipeline academico do
OpenCode e documentado em \texttt{criador-artigo/} e \texttt{basis-research/}.
49 agentes MASWOS colaboram na escrita, 12 agentes SEEKER fazem a pesquisa
bibliografica, e 5 agentes de banca simulam a revisao por pares.}:

\begin{verbatim}
SEEKER (pesquisa) → MASWOS (escrita, 49 agentes)
  → Anti-AI (TSAC, 87 palavras banidas)
  → Banca simulada (5 revisores + 4 orientadores)
  → Correcao iterativa (loop ate score >= 95/100)
  → Cora-Debate V1-V7 (verificacao simbolica)
  → Exportacao LaTeX/PDF
  → AutoEvolve (aprende com o ciclo e gera novas skills)
\end{verbatim}

\subsection{Gerando um Artigo sobre GAT com o Pipeline}

\begin{lstlisting}[caption={Gerar artigo GAT via pipeline OpenCode}]
# Terminal: gerar artigo sobre um experimento GAT
opencode run /artigo --tema "Geometric Arbitrage Theory:
  Validacao empirica do Teorema Central em mercados
  simulados com 3 ativos de Ito" --formato abnt --paginas 15

# O pipeline:
# 1. SEEKER busca referencias (arXiv, PubMed, OpenAlex)
# 2. MASWOS (49 agentes) escreve o artigo
# 3. TSAC verifica citacoes (87 palavras banidas)
# 4. Banca (5 revisores) avalia
# 5. Cora-Debate (V1-V7) audita cada afirmacao
# 6. Loop ate score >= 95/100
# 7. Exporta PDF (ABNT)
\end{lstlisting}

\section{Exercicios Praticos}

\begin{exer}[Primeiro Experimento GAT]
Instale o OpenCode Ecosystem e execute o Experimento 1 (mercado de 2 ativos).
Responda:
\begin{enumerate}[(a)]
  \item Por que $\mu \neq r$ nao constitui arbitragem?
  \item O que precisa ser verdade para que exista arbitragem?
  \item Simule 1.000 trajetorias e verifique a distribuicao de $R$ — ela e
        significativamente diferente de zero?
\end{enumerate}
\end{exer}

\begin{exer}[Calibracao dos Verificadores]
Usando o Cora-Debate, submeta 5 afirmacoes \textbf{deliberadamente falsas}
sobre um mercado simulado. Registre quais verificadores detectaram cada
erro e quais nao detectaram. O que isso revela sobre as limitacoes da
verificacao simbolica?
\end{exer}

\begin{exer}[Ciclo SDD+TDD Completo]
Escolha um aspecto da GAT (ex: condicao de Novikov, equacao de continuidade,
Teorema de Noether) e implemente o ciclo SPEC→TDD→IMPLEMENTAR→VERIFICAR→
AUDITAR. Documente o processo. Submeta ao Cora-Debate. Publique o resultado.
\end{exer}

\section{Referencias do Capitulo}

\begin{enumerate}[nosep]
  \item Farinelli, S. (2021). Geometric Arbitrage Theory and Market Dynamics
        Reloaded. arXiv:0910.1671v10.
  \item Claro, M. (2026). OpenCode Ecosystem v4.6.1. GitHub.
  \item Beck, K. (2003). Test-Driven Development: By Example. Addison-Wesley.
        ISBN: 978-0321146533.
  \item Wang, L. et al. (2024). A survey on LLM-based autonomous agents.
        Frontiers of Computer Science. DOI: 10.1007/s11704-024-40231-1.
\end{enumerate}
